% GENERATED FILE. Edit canonical lessons, then run npm run generate:books.
% canonical-source-hash: fnv1a64:91794b8eec53b65b
% canonical-lessons: LA-C531-atriensis, LA-C531-ianitor, LA-C531-ostiarius, LA-C531-lecticarius, LA-C531-vilica

\chapter{Jobs in a great house}
\label{ch:la-jobs-in-a-great-house}
\hlchaptermodality{531}

\begin{chapteropening}
\textbf{By the end of this chapter:} \emph{I can name the jobs of slaves in a large household.}

Complete the last lesson of chapter 531: I can name the jobs of slaves in a large household.
\end{chapteropening}

\section[ātriēnsis, ātriēnsis]{ātriēnsis, ātriēnsis — a steward, house-manager}

\label{lesson:LA-C531-atriensis}



\begin{quote}
Before the new one: say the Latin for of a slave, then the Latin for a footman.
\end{quote}

\subsection*{You'll want to know: ātriēnsis, ātriēnsis}
\textbf{ātriēnsis, ātriēnsis} — \textquotedblleft{}a steward, house-manager\textquotedblright{}.

\begin{quote}
Ātriēnsis servīs labōrem dat.
\end{quote}

\begin{quote}
\emph{The steward gives the slaves their work.}
\end{quote}

Masculine, third declension, from \emph{ātrium}: the slave in charge of the hall and the house.

\begin{cousinweb}
\textbf{ātriēnsis} is built on \emph{ātrium} \textquotedblleft{}hall\textquotedblright{} with \emph{-ēnsis} (Lewis \& Short; EtymDB).
\end{cousinweb}

\begin{grammarlens}[title={What you've built}]
One more word for telling, explaining and describing at B1.
\end{grammarlens}

\subsection*{Your turn}
\begin{itemize}
\raggedright
  \item \emph{Say it:} \emph{ātriēnsis, ātriēnsis}
  \item \emph{Say it:} \emph{ātriēnsis, ātriēnsis}, once more
  \item \emph{Say it:} a sentence of your own with \emph{ātriēnsis, ātriēnsis}
  \item \emph{Recall:} say the Latin for of a slave, then the Latin for a footman, then say \emph{ātriēnsis, ātriēnsis} again
\end{itemize}

\begin{tcolorbox}[breakable,colback=teal!4,colframe=teal!35!black,title={Before you move on}]
What does ātriēnsis, ātriēnsis mean? (\textquotedblleft{}A steward, house-manager\textquotedblright{}.) Say it once more.
\end{tcolorbox}
\section[iānitor, iānitōris]{iānitor, iānitōris — a doorkeeper}

\label{lesson:LA-C531-ianitor}



\begin{quote}
Before the new one: say the Latin for a footman, then the Latin for a steward.
\end{quote}

\subsection*{You'll want to know: iānitor, iānitōris}
\textbf{iānitor, iānitōris} — \textquotedblleft{}a doorkeeper\textquotedblright{}.

\begin{quote}
Iānitor ad iānuam sedet et hospitēs salūtat.
\end{quote}

\begin{quote}
\emph{The doorkeeper sits by the door and greets the guests.}
\end{quote}

Masculine, third declension, from \emph{iānua} — door. English \emph{janitor}.

\begin{cousinweb}
\textbf{iānitor} is built on \emph{iānua} \textquotedblleft{}door\textquotedblright{} with \emph{-tor} (Lewis \& Short; EtymDB). English \emph{janitor} comes from it.
\end{cousinweb}

\begin{grammarlens}[title={What you've built}]
One more word for telling, explaining and describing at B1.
\end{grammarlens}

\subsection*{Your turn}
\begin{itemize}
\raggedright
  \item \emph{Say it:} \emph{iānitor, iānitōris}
  \item \emph{Say it:} \emph{iānitor, iānitōris}, once more
  \item \emph{Say it:} a sentence of your own with \emph{iānitor, iānitōris}
  \item \emph{Recall:} say the Latin for a footman, then the Latin for a steward, then say \emph{iānitor, iānitōris} again
\end{itemize}

\begin{tcolorbox}[breakable,colback=teal!4,colframe=teal!35!black,title={Before you move on}]
What does iānitor, iānitōris mean? (\textquotedblleft{}A doorkeeper\textquotedblright{}.) Say it once more.
\end{tcolorbox}
\section[ostiārius, ostiāriī]{ostiārius, ostiāriī — a porter, doorman}

\label{lesson:LA-C531-ostiarius}



\begin{quote}
Before the new one: say the Latin for a steward, then the Latin for a doorkeeper.
\end{quote}

\subsection*{You'll want to know: ostiārius, ostiāriī}
\textbf{ostiārius, ostiāriī} — \textquotedblleft{}a porter, doorman\textquotedblright{}.

\begin{quote}
Ostiārius canem ad portam habet.
\end{quote}

\begin{quote}
\emph{The porter keeps a dog at the gate.}
\end{quote}

Masculine, second declension, from \emph{ostium} — door, entrance.

\begin{cousinweb}
\textbf{ostiārius} is built on \emph{ostium} \textquotedblleft{}door, entrance\textquotedblright{} with \emph{-ārius} (Lewis \& Short; EtymDB), and EtymDB relates \emph{ostium} to \emph{ōs} \textquotedblleft{}mouth\textquotedblright{}.
\end{cousinweb}

\begin{grammarlens}[title={What you've built}]
One more word for telling, explaining and describing at B1.
\end{grammarlens}

\subsection*{Your turn}
\begin{itemize}
\raggedright
  \item \emph{Say it:} \emph{ostiārius, ostiāriī}
  \item \emph{Say it:} \emph{ostiārius, ostiāriī}, once more
  \item \emph{Say it:} a sentence of your own with \emph{ostiārius, ostiāriī}
  \item \emph{Recall:} say the Latin for a steward, then the Latin for a doorkeeper, then say \emph{ostiārius, ostiāriī} again
\end{itemize}

\begin{tcolorbox}[breakable,colback=teal!4,colframe=teal!35!black,title={Before you move on}]
What does ostiārius, ostiāriī mean? (\textquotedblleft{}A porter, doorman\textquotedblright{}.) Say it once more.
\end{tcolorbox}
\section[lectīcārius, lectīcāriī]{lectīcārius, lectīcāriī — a litter-bearer}

\label{lesson:LA-C531-lecticarius}



\begin{quote}
Before the new one: say the Latin for a doorkeeper, then the Latin for a porter.
\end{quote}

\subsection*{You'll want to know: lectīcārius, lectīcāriī}
\textbf{lectīcārius, lectīcāriī} — \textquotedblleft{}a litter-bearer\textquotedblright{}.

\begin{quote}
Sex lectīcāriī dominam per urbem portant.
\end{quote}

\begin{quote}
\emph{Six litter-bearers carry their mistress through the city.}
\end{quote}

Masculine, second declension, from \emph{lectīca} — a litter.

\begin{cousinweb}
\textbf{lectīcārius} is built on \emph{lectīca} \textquotedblleft{}litter, sedan chair\textquotedblright{} with \emph{-ārius} (Lewis \& Short); \emph{lectīca} itself comes from \emph{lectus} \textquotedblleft{}couch\textquotedblright{} (EtymDB).
\end{cousinweb}

\begin{grammarlens}[title={What you've built}]
One more word for telling, explaining and describing at B1.
\end{grammarlens}

\subsection*{Your turn}
\begin{itemize}
\raggedright
  \item \emph{Say it:} \emph{lectīcārius, lectīcāriī}
  \item \emph{Say it:} \emph{lectīcārius, lectīcāriī}, once more
  \item \emph{Say it:} a sentence of your own with \emph{lectīcārius, lectīcāriī}
  \item \emph{Recall:} say the Latin for a doorkeeper, then the Latin for a porter, then say \emph{lectīcārius, lectīcāriī} again
\end{itemize}

\begin{tcolorbox}[breakable,colback=teal!4,colframe=teal!35!black,title={Before you move on}]
What does lectīcārius, lectīcāriī mean? (\textquotedblleft{}A litter-bearer\textquotedblright{}.) Say it once more.
\end{tcolorbox}
\section[vīlica, vīlicae]{vīlica, vīlicae — a farm housekeeper}

\label{lesson:LA-C531-vilica}



\begin{quote}
Before the new one: say the Latin for a porter, then the Latin for a litter-bearer.
\end{quote}

\subsection*{You'll want to know: vīlica, vīlicae}
\textbf{vīlica, vīlicae} — \textquotedblleft{}a farm housekeeper\textquotedblright{}.

\begin{quote}
Vīlica cēnam familiae rūsticae parat.
\end{quote}

\begin{quote}
\emph{The farm housekeeper prepares dinner for the farm slaves.}
\end{quote}

Feminine, first declension: the wife of the \emph{vīlicus}, the farm manager. She ran the farm kitchen and stores.

\begin{cousinweb}
\textbf{vīlica} is the feminine of \emph{vīlicus} \textquotedblleft{}farm manager\textquotedblright{}, built on \emph{vīlla} \textquotedblleft{}country house, farm\textquotedblright{} (Lewis \& Short).
\end{cousinweb}

\begin{grammarlens}[title={What you've built}]
One more word for telling, explaining and describing at B1.
\end{grammarlens}

\subsection*{Your turn}
\begin{itemize}
\raggedright
  \item \emph{Say it:} \emph{vīlica, vīlicae}
  \item \emph{Say it:} \emph{vīlica, vīlicae}, once more
  \item \emph{Say it:} a sentence of your own with \emph{vīlica, vīlicae}
  \item \emph{Recall:} say the Latin for a porter, then the Latin for a litter-bearer, then say \emph{vīlica, vīlicae} again
\end{itemize}

\begin{tcolorbox}[breakable,colback=teal!4,colframe=teal!35!black,title={Before you move on}]
What does vīlica, vīlicae mean? (\textquotedblleft{}A farm housekeeper\textquotedblright{}.) Say it once more.
\end{tcolorbox}
